\section{Introduction and setup}

\subsection{Introduction}
% TODO (Member 1): problem statement, the spam-filter application, why
% accuracy is the wrong target, and a roadmap of Sections 2-7.
\emph{To be written by Member 1.}

\subsection{Experimental setup}
% Member 6 (C1)

\paragraph{Application, actions and losses.}
We filter email. Label $y=1$ means spam and $y=0$ a real email. For each email the filter
chooses one of three actions: deliver to the inbox ($a=0$), send to the spam folder ($a=1$), or
hold for human inspection ($a=r$; coded as \texttt{2}). Correct decisions cost nothing. Sending
a real email to spam costs $\cfp$, delivering spam costs $\cfn$, and inspection costs $\ci$
whatever the label. We assume the inspector always finds the true label and routes the email
correctly, so no further loss follows an inspection. Unless stated otherwise $\cfp=10$ and $\cfn=1$;
all experiments sweep $\ci$ or the ratio $\cfp/\cfn$.

\paragraph{Data.}
\emph{Spambase}~\cite{spambase} has $4{,}601$ emails, of which $1{,}813$ ($39.4\%$) are spam, and
57 non-negative features: 48 word frequencies, 6 character frequencies and 3 run-length
statistics of capital letters. The features are heavily skewed and $77\%$ of all feature values are
exactly zero, which matters for the Gaussian models (Section~\ref{sec:models}).
The \emph{synthetic} data are two-class Gaussian samples with chosen means, covariances and
prior $\pio$. Because $p(x\mid y)$ is known we can compute the exact posterior and Bayes risk,
and change the class proportions exactly (Section~6). Section~7 also uses a one-dimensional
special case with closed-form answers.

\paragraph{Splits and error bars.}
Each data set is split into stratified $60/20/20$ train, validation and test parts
($2{,}761/921/919$ emails for Spambase, each with the overall spam rate). Models are fitted on the
training part. Every choice (regularisation strength, temperature, anything else) is made on the
validation part. The test part is used only to report results. We repeat everything over 20 splits
(seeds $0$--$19$) and report the mean with a standard-error bar.

\paragraph{Preprocessing.}
Standardisation (zero mean, unit variance), optionally after the transform $\log(1+x)$, uses
statistics of the training part only and is then applied unchanged to validation and test data.

\paragraph{Prior shift.}
To simulate a deployment prior $\pio'$ we resample test emails \emph{within each class}, so
$p(x\mid y)$ is unchanged and only the label proportions move. By default we sample without
replacement and keep the largest set the data allow.

\paragraph{Code and reproducibility.}
All models are written from scratch in Python with NumPy; scikit-learn is used only to cross-check
numbers. The shared modules \texttt{data.py}, \texttt{decision.py} and \texttt{metrics.py} implement
the interface above. A unit test checks the decision rule of Section~2 against a brute-force
minimisation of the posterior expected loss on $10{,}000$ random posteriors and loss matrices,
covering cases both with and without an inspect band. The command
\texttt{python -m project4.run\_all} regenerates every figure and table in this report from scratch.
